\section{Conclusion}

We have presented a simple data-constrained scaling law that models the cost of data repetition with a simple, additive overfitting penalty. A complexity ladder of one-, two-, and four-parameter forms provides a Pareto frontier of fit quality versus complexity, with even the simplest form substantially outperforming prior effective-data formulations.
The overfitting penalty provides a new axis for evaluating training configurations, directly quantifying robustness to data repetition. This property is of increasing importance as compute grows faster than data. As a case study, we show that strong weight decay ($\lambda=1.0$) reduces the overfitting penalty by approximately 70\%, providing prescriptive guidance that explains, and can replace, the expensive per-configuration hyperparameter sweeps required by prior empirical work~\citep{kim2026pretraining}.
Finally, our overfitting penalty yields qualitatively new compute-optimal allocation advice: beyond a threshold compute budget, further data repetition is counterproductive, and practitioners should scale model capacity instead. Together, these findings provide concrete guidance for the increasingly common setting where high-quality data, not compute, is the binding constraint.
